\documentclass[UTF8,a4paper,12pt]{ctexart}
\usepackage{amsmath, amssymb}
\usepackage{graphicx}
\usepackage{booktabs}
\usepackage{float}
\usepackage[sort&compress]{gbt7714}
\bibliographystyle{gbt7714-numerical}
\usepackage[hidelinks]{hyperref}

\title{论文题目}
\author{作者\thanks{单位，邮箱}}
\date{\today}

\begin{document}
\maketitle

\begin{abstract}
在这里写摘要：研究的问题、方法、主要结果和结论。
\end{abstract}

\noindent\textbf{关键词：}关键词一；关键词二；关键词三

\section{引言}

说明研究背景和问题。引用文献写作 \cite{lecun2015deep}，多篇文献写作 \cite{lecun2015deep,gbt7714-2015}。

\section{方法}

公式用 equation 环境，并给出标签以便引用：
\begin{equation}
  y = f(\mathbf{x}; \theta) + \varepsilon
  \label{eq:model}
\end{equation}
式~\eqref{eq:model} 中 $\theta$ 为待估参数。

\section{结果}

表格用 booktabs 三线表：
\begin{table}[H]
  \centering
  \caption{实验结果}
  \label{tab:results}
  \begin{tabular}{lcc}
    \toprule
    方法 & 指标一 & 指标二 \\
    \midrule
    方法 A & 0.91 & 0.88 \\
    方法 B & 0.94 & 0.90 \\
    \bottomrule
  \end{tabular}
\end{table}

表~\ref{tab:results} 给出了结果。

\section{结论}

总结主要发现和后续工作。

\bibliography{refs}

\end{document}
